
\SGASection{6}{Functor of Cartan subgroups and functor of parabolic
subgroups}
\SGARefTarget{sga3:xv:section:6}

When \(G\) is a smooth connected algebraic group defined over an
algebraically closed field \(k\), sub-tori of \(G\), maximal
sub-tori, Cartan subgroups (Exposé XII~1), Borel subgroups (Exposé
\SGARefLink{sga3:xiv:4.1}{XIV~4.1}), and parabolic subgroups (Exposé \SGARefLink{sga3:xiv:4.8.1}{XIV~4.8 bis}) have already
been defined. We extend these notions to a group prescheme over an
arbitrary base as follows.

\textbf{\SGAInlineRefTarget{sga3:xv:definition:6-1}Definition 6.1.} \emph{Let \(S\) be a prescheme, let \(G\)
be an \(S\)-group prescheme of finite presentation, let \(H\) be an
\(S\)-group prescheme, and let \(i:H\to G\) be an
\(S\)-monomorphism making \(H\) a subgroup of \(G\). We say that
\(H\) is a sub-torus of \(G\) (respectively, a maximal sub-torus, a
Cartan subgroup, a Borel subgroup, or a parabolic subgroup) if:}

\begin{enumerate}
\renewcommand{\labelenumi}{\textup{\roman{enumi})}}
\item
  \emph{i) \(H\) is smooth over \(S\).}
\item
  \emph{ii) For every geometric point \(\bar s\) above \(S\),
  \(H_{\bar s}\) is a sub-torus of
  \((G_{\bar s})^{0}_{\mathrm{red}}\) (respectively, a maximal
  sub-torus, a Cartan subgroup, a Borel subgroup, or a parabolic
  subgroup).}
\end{enumerate}
\SGARefTarget{sga3:xv:definition:6:1}

\textbf{\SGAInlineRefTarget{sga3:xv:remark:6-1-bis}Remarks 6.1 bis.}

\begin{enumerate}
\renewcommand{\labelenumi}{\textup{\alph{enumi})}}
\item
  If the \(S\)-group \(H\) is a sub-torus of \(G\)
  (respectively, \(\ldots\)), its fibres are connected, and
  consequently \(H\) is of finite presentation over \(S\) (Exposé
  VI\textsubscript{B}~\SGARefLink{par:sga3-VIA-5.3.3}{5.3.3}).
\item
  If \(H\) is a sub-torus of \(G\), then \(H\) is a torus in the
  sense of Exposé IX, as follows immediately from Exposé \SGARefLink{sga3:x:theorem:8-1}{X~8.1}.
  Moreover, the monomorphism \(i:H\to G\) is an immersion (see \SGARefLink{sga3:xv:lemma:8:3}{8.3}
  below).
\item
  If \(G\) is smooth over \(S\), with connected fibres, and if
  \(H\) is a Cartan subgroup of \(G\) (respectively, a Borel
  subgroup or a parabolic subgroup), the monomorphism \(i:H\to G\)
  is an immersion. Thus our definitions coincide with those
  introduced in Exposés XII and XIV. Indeed, \(H\) is then equal to
  its connected normalizer (by \SGARefLink{sga3:xii:theorem:6-6:item:c}{XII~6.6(c)}, \SGARefLink{sga3:xiv:4.8}{XIV~4.8} and \SGARefLink{sga3:xiv:paragraph:4-8-bis}{4.8 bis}),
  and it suffices to apply \SGARefLink{sga3:xv:remark:5:8}{5.8}.
\end{enumerate}

\textbf{Definition 6.1 ter.} \emph{Let \(S\) be a prescheme, let
\(G\) be an \(S\)-group prescheme locally of finite type, and let
\(s\) be a point of \(S\). The \textbf{nilpotent rank} of \(G\) at
\(s\), denoted by \(\rho_n(s)\), is the dimension of the Cartan
subgroups of \((G_s)^0_{\mathrm{red}}\). The reductive rank
\(\rho_r(s)\), unipotent rank \(\rho_u(s)\), and abelian rank
\(\rho_{\mathrm{ab}}(s)\) are defined similarly (see Exposé
\SGARefLink{sga3:x:remark:8-7}{X~8.7}).}
\SGARefTarget{sga3:xv:definition:6:1:ter}

If \(G\) is now a smooth connected algebraic group defined over an
algebraically closed field \(k\), recall that the \textbf{radical}
of \(G\), denoted by \(\operatorname{rad}(G)\), is the largest
invariant smooth connected solvable algebraic subgroup of \(G\);
\(G/\operatorname{rad}(G)\) is then semisimple (use Exposé \SGARefLink{sga3:xii:6.1}{XII~6.1}
to reduce to the affine case). If \(G\) is moreover affine, the
\textbf{unipotent radical} \(\operatorname{rad}^{u}(G)\) is defined
as the largest invariant smooth connected unipotent algebraic
subgroup of \(G\); \(G/\operatorname{rad}^{u}(G)\) is then
reductive.


\textbf{\SGAInlineRefTarget{sga3:xv:proposition:6-2}Proposition 6.2.} \emph{Let \(S\) be a prescheme, let \(G\)
and \(H\) be two \(S\)-group preschemes of finite presentation, and
let \(i:H\to G\) be an \(S\)-monomorphism making \(H\) a subgroup
of \(G\). Let \(P(s)\) be one of the following properties concerning
a point \(s\) of \(S\):}


\begin{enumerate}
\renewcommand{\labelenumi}{\textup{\roman{enumi})}}
\item \emph{\((G_{\bar s})^{0}_{\mathrm{red}}\) is an abelian
variety (respectively, is affine, is a torus, or is unipotent).}
\item \emph{\((H_{\bar s})^{0}_{\mathrm{red}}\) is a maximal torus
of \(G_{\bar s}\).}
\item\SGAInlineRefTarget{sga3:xv:proposition:6-2:item:iii} \emph{\((H_{\bar s})^{0}_{\mathrm{red}}\) is the centralizer
in \((G_{\bar s})^{0}_{\mathrm{red}}\) of a torus of \(G_{\bar s}\)
(respectively, is a Cartan subgroup of \(G_{\bar s}\)).}
\item \emph{\((H_{\bar s})^{0}_{\mathrm{red}}\) is a Borel subgroup
(respectively, a parabolic subgroup) of \(G_{\bar s}\).}
\item \emph{\((H_{\bar s})^{0}_{\mathrm{red}}\) is the radical of
\(G_{\bar s}\) (respectively,
\((G_{\bar s})^{0}_{\mathrm{red}}\) is semisimple).}
\item \emph{\(G_{\bar s}\) is affine and
\((H_{\bar s})^{0}_{\mathrm{red}}\) is the unipotent radical of
\(G_{\bar s}\) (respectively,
\((G_{\bar s})^{0}_{\mathrm{red}}\) is reductive).}
\end{enumerate}

\emph{Then the set \(E\) of points \(s\) of \(S\) for which \(P(s)\)
holds is locally constructible (EGA \(0_{\mathrm{III}}\), 9.1.11).}
\SGARefTarget{sga3:xv:proposition:6:2}

\textbf{\SGAInlineRefTarget{sga3:xv:remark:6-2-1}Remark 6.2.1.} This proposition complements Exposé
VI\textsubscript{B}, \S10. The structure of \(E\) can be described
more precisely by using semicontinuity theorems (see Exposé \SGARefLink{sga3:x:remark:8-7}{X~8.7});
we shall see an example below.

\SGAProof{Proof of \SGARefLink{sga3:xv:proposition:6-2}{6.2}}
Note that if \(S\) is the spectrum of a field, \(E\) is invariant
under extension of that field. A standard reduction (EGA IV~9)
then allows us to reduce to the case where \(S\) is noetherian and
integral, with generic point \(\eta\). One must show that \(E\) or
\(\operatorname{ens}(S)\setminus E\) contains a neighbourhood of
\(\eta\) (EGA IV~9.2.1). One
may assume \(S\) affine of ring \(A\) and field of fractions \(K\). If \(L\) is a finite extension of \(K\), it is immediate that

there exists an \(A\)-subalgebra \(B\) of \(L\), finite over \(A\), having \(L\) as field of fractions. The canonical morphism \(S'

\to S\), where \(S' = \operatorname{Spec} B\), is dominant, of finite presentation, so the image of a non-empty open subset

of \(S'\) contains a non-empty open subset of \(S\) (EGA IV 1.8.4). From the viewpoint that interests us, we may therefore

replace \(S\) by \(S'\), hence replace \(K\) by a finite extension \(L\). Thus we may choose \(L\) so that \((G_{L})_{red}\) and

\((H_{L})_{red}\) are smooth over \(L\) (EGA IV 4.6.6). Possibly restricting \(S'\), we may assume that \(G_{red}\) and

\(H_{\mathrm{red}}\) are group preschemes smooth over \(S\)
(Exposé VI\textsubscript{B}, \S10, and EGA IV~17). In view of the
properties to be proved, we may replace \(G\) and \(H\) by their
reduced identity components (Exposé VI\textsubscript{B}~\SGARefLink{lem:sga3-VIB-10.9}{10.9}), and
hence assume that \(G\) and \(H\) are smooth over \(S\) with
connected fibres. Finally, we may assume that \(H\) is a closed
subgroup prescheme of \(G\) (Exposé VI\textsubscript{B}~10.4).


\emph{Proof of i).} Possibly after a finite extension of \(K\), we
may assume that \(G_{\eta}\) admits a ``Chevalley decomposition'',
i.e. is an extension of an abelian variety \(D_{\eta}\) by a smooth
connected linear algebraic group \(F_{\eta}\) (Séminaire Bourbaki
1956/57, no.~145). By Exposé VI\textsubscript{B}~10.16, there is an
open neighbourhood \(U\) of \(\eta\) over which this generic
extension comes from an extension
\[
  1\longrightarrow F\longrightarrow G|_U\longrightarrow D
  \longrightarrow 1.
\]
One may moreover assume that \(F\) and \(D\) are smooth over \(U\)
with connected fibres, that \(F\) is affine over \(U\), and that
\(D\) is proper over \(U\) (EGA IV~8, 9 and 17). For every point
\(s\) of \(U\), \((D_s,F_s)\) is then the Chevalley decomposition of
\(G_s\). Consequently, \(G_s\) is an abelian variety (respectively,
is affine) if and only if \(F_s\) (respectively, \(D_s\)) is the unit
group; this is a constructible condition (EGA IV~9.2.6.1).


To establish the last two assertions of (i), we may, by the
preceding, assume that \(G\) is affine over \(S\). Let \(q\) be a
prime number invertible on \(S\), and let \({}_qG\) be the
``kernel'' of the \(q\)-th-power morphism in \(G\). It follows
easily from the structure of affine algebraic groups that \(G_s\)
is a torus (respectively, is unipotent) if and only if
\begin{enumerate}
\renewcommand{\labelenumi}{\textup{\alph{enumi})}}
\item \({}_qG_s\) is quasi-finite, which is a constructible
property (EGA IV~9.3.2), and
\item \({}_qG_s\) has \(q^r\) geometric points, where \(r\) is the
relative dimension of \(G\) over \(S\) (respectively,
\({}_qG_s\) has a single point).
\end{enumerate}
The function
\[
 s\longmapsto
 \bigl(\text{number of geometric points of }{}_qG_s\bigr)
\]
is constructible (EGA IV~9.7.9). This completes the proof of (i).


\emph{Proof of iii).} \textbf{a) Case of the centralizer of a
torus.} Suppose
\[
  H_{\eta}=\operatorname{Centr}_{G_{\eta}}(T_{\eta}),
\]
where \(T_{\eta}\) is a torus of \(G_{\eta}\). We shall show that
\(H_s\) is the centralizer of a sub-torus of \(G_s\) for \(s\) in a
neighbourhood of \(\eta\). By (i), after possibly restricting \(S\),
we may assume that \(T_{\eta}\) comes from a sub-torus \(T\) of
\(G\). Then \(Z=\operatorname{Centr}_{G}(T)\) is represented
(Exposé \SGARefLink{sga3:xi:corollary:6:11}{XI~6.11}) by a subgroup prescheme of \(G\). Since \(H\) and
\(Z\) coincide generically, they coincide over a neighbourhood of
\(\eta\). Thus the set \(E\) of points \(s\) of \(S\) for which
\(H_s\) is the centralizer in \(G_s\) of a sub-torus of \(G_s\) is
ind-constructible (EGA IV~1). This will suffice in Lemma~6.6 below
to prove that \(E\) is open in \(S\), and hence, \emph{a fortiori},
locally constructible.

\textbf{b) Case of a Cartan subgroup.} Suppose that \(H_{\eta}\) is
a Cartan subgroup of \(G_{\eta}\). We shall show that \(H_s\) is a
Cartan subgroup of \(G_s\) for every \(s\) in a neighbourhood of
\(\eta\). The group \(H_{\eta}\) is the centralizer in \(G_{\eta}\)
of a torus of \(G_{\eta}\), and is nilpotent (Exposé \SGARefLink{sga3:xii:6.6}{XII~6.6}). By
\SGARefLink{sga3:xii:6.6:item:a}{(a)} and Exposé VI\textsubscript{B}~\SGARefLink{sga3:xv:definition:8-4}{8.4}, \(H_s\) has the same
properties for every \(s\) in some neighbourhood \(U\) of \(\eta\).
For every \(s\in U\), the group \(H_s\) therefore has the same
reductive rank as \(G_s\), and its unique maximal torus is central
(Exposé \SGARefLink{sga3:xii:6.7}{XII~6.7}). Hence \(H_s\) is the centralizer of a maximal torus
of \(G_s\), that is, a Cartan subgroup of \(G_s\).


Suppose now that \(H\) is not a Cartan subgroup of \(G\), and let us show that \(H_{s}\) is not a Cartan subgroup of \(G_{s}\)

for \(s\) in a neighborhood \(U\) of \(\eta\). In view of the assertion provisionally admitted in (a) above, we may restrict

to the case where \(H\) is the centralizer in \(G\) of a sub-torus \(T\). But then \(H_{\eta}\) contains a Cartan subgroup

\(C_{\eta}\) of \(H_{\eta}\). We have just seen that, possibly restricting \(S\), \(C_{\eta}\) extends to a Cartan subgroup \(C\)

of \(G\), which one may assume contained in \(H\). By hypothesis \(H_{\eta}\) strictly contains \(C_{\eta}\), so \(H_{s}\)
strictly contains \(C_{s}\) for \(s\) in a neighborhood \(U\) of \(\eta\) (EGA IV 9.5.2); \emph{a fortiori}, \(H_{s}\) is not


a Cartan subgroup of \(G_{s}\) for \(s\) in \(U\).



\emph{Proof of ii).} Suppose \(H_{\eta}\) is a maximal torus of \(G_{\eta}\) and let \(C_{\eta}\) be its centralizer in \(G_{\eta}\).

By i) and iii), \(H\) is a torus over a neighborhood \(U\) of \(\eta\) and \(C = \operatorname{Centr}_{G|U}(H|U)\) is a Cartan subgroup of
\(G|U\). To prove that \(H\) is a maximal torus of \(G\) over a neighborhood of \(\eta\), we may then replace \(G\) by \(C\), then

by the linear component \(F\) of a Chevalley decomposition of \(C\) (cf.~i)). Let \(q\) be an integer invertible on \(S\),

\({}_{q}F\) the kernel of the \(q\)-th-power morphism in \(F\). Since \(F_{s}\) is affine, nilpotent, smooth and connected,
\(F_{s}\) is the direct product of its maximal torus \(T_{s}\) by a unipotent group (\emph{Bible} 6-04), so
\({}_{q}F_s={}_{q}T_s\). Since \(H_{\eta}\) is a maximal torus,
\({}_{q}H_{\eta}={}_{q}F_{\eta}\), and consequently
\({}_{q}H={}_{q}F\) over a neighborhood \(V\) of \(\eta\).
For every point \(s\) of \(V\), \({}_{q}H_s={}_{q}T_s\), so
\(H_{s} = T_{s}\) is a maximal torus.


Suppose now that \(H_{\eta}\) is not a maximal torus of \(G_{\eta}\). By i), we may restrict to the case where \(H_{\eta}\) is

a torus, then assume that it is contained in a strictly larger torus \(T_{\eta}\). The latter extends to a torus \(T\)

strictly containing \(H\) over a neighborhood \(U\) of \(\eta\). \emph{A fortiori}, \(H_{s}\) is not a maximal torus for \(s \in U\).


\emph{Proof of iv).} Possibly restricting \(S\), we may assume that the center \(Z\) of \(G\) is representable (Exp. VI\textsubscript{B} 10.11)
and flat over \(S\), as well as the quotient \(G/Z\) (\emph{loc. cit.}). The property ``\(H_{s}\) contains \(Z_{s}\)'' is
constructible (EGA IV 9.5.2) and every parabolic subgroup of \(G_{s}\) contains \(Z_{s}\) (Exp. XIV 4.9 a)); this allows us

to replace \(G\) by \(G/Z\), hence assume \(G\) affine over \(S\) (Exp. XII 6.1 and i)). We may further



assume that \(G/H\) is representable, but then \(H_{s}\) is a parabolic subgroup of \(G_{s}\) if and only if \((G/H)_{s}\) is

proper (\emph{Bible} 6 Th. 4 b)), which is an ind-constructible property (EGA IV 9.3.5). So \(E\) is ind-constructible, and

this will suffice for us to prove that \(E\) is open (Lemma 6.6), hence locally constructible.



Let us now examine the case of Borel subgroups. If \(H_{\eta}\) is a Borel subgroup of \(G_{\eta}\), i.e.~a solvable

parabolic subgroup of \(G_{\eta}\), what precedes and Exposé VI\textsubscript{B}, \SGARefLink{prop:sga3-VIB-8.4}{8.4}, entail that these properties remain true at every
point \(s\) of a neighborhood of \(\eta\). If now \(H_{\eta}\) is not a Borel subgroup of \(G_{\eta}\), to prove that the same

holds at points \(s\) of a neighborhood of \(\eta\), we may restrict (in view of what precedes) to the case where \(H_{\eta}\)

is a parabolic subgroup, then assume that \(H_{\eta}\) contains a Borel subgroup \(B_{\eta}\). We have just shown that the

latter extends to a Borel subgroup \(B\) of \(G\) over a neighborhood \(U\) of \(\eta\). Since \(H_{\eta}\) strictly contains
\(B_{\eta}\), \(H_{s}\) strictly contains \(B_{s}\) at every point of an open subset \(V\), and \(H_{s}\) is not a Borel subgroup
of \(G_{s}\) for \(s \in V\).



\emph{Proof of v).} Suppose \(H_{\eta}\) is the radical of \(G_{\eta}\). The group \(H_{\eta}\) is then invariant in \(G_{\eta}\),

solvable (smooth and connected); the same is therefore true for \(H_{s}\) for \(s\) belonging to a neighborhood \(U\) of

\(\eta\) (Exp. VI\textsubscript{B} 8 and 10), so for \(s \in U\), \(H_{s}\) is contained in the radical of \(G_{s}\). Replacing \(G\) by \(G/H\)
(Exp. VI\textsubscript{B} 10), we must


prove that if \(G_{\eta}\) is semisimple, then \(G_{s}\) is semisimple at every point of a neighborhood \(V\) of \(\eta\). Using

i) and ii), one may assume that \(G\) is affine over \(S\) and possesses a maximal torus \(T\). Let \(W\) be the Weyl group of

\(T\) (\SGARefLink{sga3:xii:section:2}{Exp. XII 2}), which is quasi-finite and étale over \(S\), hence finite and étale over an open subset \(V\). It then

follows from the elementary properties of roots (\SGARefLink{sga3:xix:proposition:1-12}{Exp. XIX 1.12}) that \(G\) is semisimple over \(V\).



Suppose now that \(H_{\eta}\) is not the radical of \(G_{\eta}\). Possibly replacing \(K\) by a finite extension \(L\), we may

assume that \(G_{\eta}\) admits a radical \(R_{\eta}\). By what precedes, \(R_{\eta}\) extends to a subgroup prescheme \(R\) of

\(G|_{U}\) such that for every \(s \in U\), \(R_{s}\) is the radical of \(G_{s}\). By hypothesis, \(R_{\eta} \neq H_{\eta}\). So

\(R_{s} \neq H_{s}\) for \(s \in V\). It remains to prove that if \(G_{\eta}\) is not semisimple, neither is \(G_{s}\) at

neighboring points, but this is a particular case of what precedes (take \(H\) = unit group).



\emph{Proof of vi).} The proof is entirely analogous to that of v), in view of i), and is left to the care of the reader.



\textbf{\SGAInlineRefTarget{sga3:xv:corollary:6-3}Corollary 6.3.} \emph{Let \(S_0\) be a quasi-compact
prescheme, let \((S_i)_{i\in L}\) be a projective system of
\(S_0\)-preschemes affine over \(S_0\), and let
\(S=\varprojlim S_i\) (EGA IV~8.2). Let \(G_0\) be a group
prescheme of finite presentation over \(S_0\), set
\[
  G_i=G_0\times_{S_0}S_i,\qquad G=G_0\times_{S_0}S,
\]
and let \(H\) be a subgroup of \(G\). If \(H\) is a sub-torus of
\(G\) (respectively, a maximal sub-torus, a Cartan subgroup, a Borel
subgroup, or a parabolic subgroup), then there exist an index
\(i\in L\) and a subgroup \(H_i\) of \(G_i\) such that
\[
  H=H_i\times_{S_i}S
\]
and \(H_i\) is a sub-torus of \(G_i\) (respectively, \(\ldots\)).}
\SGARefTarget{sga3:xv:corollary:6:3}

Indeed, \(H\) is smooth with connected fibres, hence of finite
presentation over \(S\) (Exposé VI\textsubscript{B}~\SGARefLink{par:sga3-VIA-5.3.3}{5.3.3}). By
Exposé VI\textsubscript{B}, \S10, there exist \(i\in L\) and a
subgroup \(H_i\) of \(G_i\), smooth over \(S_i\), such that
\(H=H_i\times_{S_i}S\). Corollary~6.3 then follows from
Definition~6.1, Proposition~6.2 and EGA IV~9.3.3.


\textbf{Corollary 6.3 bis.} \emph{Let \(S\) be a prescheme and let
\(G\) be an \(S\)-group prescheme of finite presentation. Then the
functions \(\rho_n,\rho_r,\rho_u,\rho_{\mathrm{ab}}\) (cf.~\SGARefLink{sga3:xv:definition:6:1:ter}{6.1 ter})
are locally constructible on \(S\).}
\SGARefTarget{sga3:xv:corollary:6:3:bis}


It suffices to show (EGA IV 9.) that if \(S\) is a noetherian integral scheme with generic point \(\eta\), the functions in

question are constant on a neighborhood of \(\eta\). Possibly replacing \(S\) by a scheme \(S'\) finite over \(S\) and

dominating \(S\), we may assume that \(G_{\eta}\) admits a Cartan subgroup \(C_{\eta}\) with a Chevalley decomposition \(1 \to

L_{\eta} \to C_{\eta} \to A_{\eta} \to 1\). The argument made in \SGARefLink{sga3:xv:proposition:6-2}{6.2} i) proves that this decomposition extends to a

Chevalley decomposition over a neighborhood of \(\eta\):



\[ 1 \longrightarrow L \longrightarrow C \longrightarrow A \longrightarrow 1. \]



Moreover, one may assume that \(C\) is a Cartan subgroup of \(G\) (\SGARefLink{sga3:xv:corollary:6-3}{6.3}) and that the maximal torus \(T_{\eta}\) of \(L_{\eta}\)

extends to a maximal torus \(T\) of \(L\) (\SGARefLink{sga3:xv:corollary:6-3}{6.3}). The corollary follows immediately from this and from the definitions.



\SGARefTarget{sga3:xv:subsection:6:4:0}
\textbf{\SGAInlineRefTarget{sga3:xv:paragraph:6-4-0}6.4.0.} Let \(S\) be a prescheme, \(G\) an \(S\)-prescheme in groups of finite
presentation, and \(\mathcal{L}\) the functor of subgroups of \(G\) that are smooth,
have connected fibres, and are equal to their connected normalizer (cf.~§5).
The functor \(\mathcal{L}\) is representable (\SGARefLink{sga3:xv:theorem:5-2}{5.2} and \SGARefLink{sga3:xv:remark:5:8}{5.8}). The rest of this section
is devoted to certain subfunctors of \(\mathcal{L}\). More precisely, define
\(\mathcal{L}_{C}\) (resp. \(\mathcal{CT}\), resp. \(\mathcal{P}\)) as follows:
for every \(S\)-prescheme \(S'\), the set \(\mathcal{L}_{C}(S')\)
(resp. \(\mathcal{CT}(S')\), resp. \(\mathcal{P}(S')\)) consists of the subgroups
\(H\) of \(G_{S'}\) that are smooth over \(S'\), have connected fibres, are equal
to their connected normalizer, and are such that, for every point \(s'\) of \(S'\),
\(H_{s'}\) contains a Cartan subgroup (\SGARefLink{sga3:xv:definition:6-1}{6.1}) of \(G_{s'}\)
(resp. is the centralizer in \((G_{s'})^{0}_{\mathrm{red}}\) of a subtorus of
\(G_{s'}\), resp. is a parabolic subgroup (\SGARefLink{sga3:xv:definition:6-1}{6.1}) of \(G_{s'}\)).

\SGARefTarget{sga3:xv:theorem:6:4}
\textbf{Theorem 6.4.} \emph{Let \(S\) be a prescheme and \(G\) an \(S\)-prescheme in
groups of finite presentation. The \(S\)-functors \(\mathcal{L}_{C}\),
\(\mathcal{CT}\), and \(\mathcal{P}\) defined in \SGARefLink{sga3:xv:paragraph:6-4-0}{6.4.0} are representable by
\(S\)-preschemes that are quasi-projective and of finite presentation over \(S\).}

\SGARefTarget{sga3:xv:remark:6:5}
\textbf{Remark 6.5.} If \(G\) has smooth fibres---for example, if \(G\) is smooth
over \(S\), or if the residue characteristics of \(S\) are zero
(Exposé VI\textsubscript{B}, \SGARefLink{cor:sga3-VIB-1.6.1}{1.6.1})---then every subgroup \(H\) of \(G\), smooth over
\(S\) with connected fibres, such that \(H_s\) contains a Cartan subgroup of
\((G_s)^0\) for every point \(s\) of \(S\), is necessarily equal to its connected
normalizer (and hence is an element of \(\mathcal{L}_{C}(S)\)). Indeed, by \SGARefLink{sga3:xv:proposition:5-1}{5.1} iii)
we may assume that \(S\) is the spectrum of a field, in which case the property was
noted at the end of the statement of \SGARefLink{sga3:xiii:theorem:2:1}{Exposé XIII, 2.1}.

There are natural monomorphisms
\[
\begin{tikzcd}[column sep=huge,row sep=large]
\mathcal{CT} \arrow[dr,hook] & \\
& \mathcal{L}_{C} \arrow[r,hook] & \mathcal{L}.\\
\mathcal{P} \arrow[ur,hook] &
\end{tikzcd}
\]


Let us show that these monomorphisms are open immersions of finite presentation (which will already prove, in view of

\SGARefLink{sga3:xv:theorem:5-2}{5.2}, that \(\mathcal{L}_{C}\), \(\mathcal{CT}\), \(\mathcal{P}\) are representable by \(S\)-preschemes, union of an increasing

sequence of open subpreschemes, quasi-projective and of finite presentation over \(S\)). Now this will follow, by the

usual technique, from the following lemma:



\SGARefTarget{sga3:xv:lemma:6:6}
\textbf{Lemma 6.6.} \emph{Let \(S\) be a prescheme, \(G\) an \(S\)-prescheme in
groups of finite presentation, and \(H\) a subgroup of \(G\), smooth with connected
fibres. The set of points \(s\) of \(S\) such that \(H_s\) contains a Cartan subgroup
of \(G_s\) (resp. is equal to the centralizer in \((G_s)^0_{\mathrm{red}}\) of a torus
of \(G_s\)) is open in \(S\). If, moreover, \(H\) is equal to its connected
normalizer,\footnote{This hypothesis is in fact superfluous; see Exposé XVII,
Appendix III, 3.} then the set of points \(s\) of \(S\) such that \(H_s\) is a
parabolic subgroup of \(G_s\) is also open in \(S\).}


The assertion to be proved is local on \(S\), which allows us to assume \(S\) affine, then, by EGA IV 8.1 and 6.3, \(S\)

noetherian. Denote by \(E\) the set of points of \(S\) having the property in question. It then follows from the assertions

effectively proved in 6.2 that \(E\) is ind-constructible. But \(S\) is noetherian, so to prove that \(E\) is open, it

suffices to show that \(E\) is stable under generizations (EGA IV 1.10.1). Using EGA II 7.1.9, we are finally led to prove

that if \(S\) is the spectrum of a discrete valuation ring, and if the closed point \(s\) belongs to \(E\), then so does the

generic point \(t\).



We shall need the following lemma:



\SGARefTarget{sga3:xv:lemma:6:7}
\textbf{Lemma 6.7.} \emph{Let \(A\) be a complete noetherian local ring,
\(S=\operatorname{Spec}A\), \(s\) the closed point of \(S\), \(H\) an
\(S\)-prescheme in groups, smooth with connected fibres, and \(T_s\) a subtorus of
\(H_s\). Then:
\begin{enumerate}
\renewcommand{\labelenumi}{\roman{enumi})}
\item there exists a closed subgroup prescheme \(C\) of \(H\), smooth with connected
fibres, such that \(C_s=\operatorname{Centr}_{H_s}(T_s)\);
\item for every point \(t\) of \(S\), \(C_t\) is the centralizer in \(H_t\) of a
subtorus \(T_t\) of \(H_t\).
\end{enumerate}}

\emph{Proof of \SGARefLink{sga3:xv:lemma:6:7}{6.7}.} Let \(T'\) be an \(S\)-torus for which there is an isomorphism
\(u_0:T'_s\xrightarrow{\sim}T_s\) (\SGARefLink{sga3:x:corollary:4-6}{Exposé X, 4.6}). Let \(\mathfrak m\) be the
maximal ideal of \(A\), and put \(A_n=A/\mathfrak m^n\),
\(S_n=\operatorname{Spec}A_n\), \(H_n=H\times_S S_n\), and so on. Since \(H\) is
smooth over \(S\), for every integer \(n>0\) there is an \(S_n\)-group morphism
\[
u_n:T'_n\longrightarrow H_n
\]
lifting \(u_0\) (\SGARefLink{sga3:ix:theorem:3:6}{Exposé IX, 3.6}), and by induction on \(n\) we may assume that
\(u_n\) lifts \(u_{n-1}\). Let \(q\) be a prime invertible on \(S\). For every
integer \(\ell\) that is a power of \(q\), write \({}_{\ell}u_n\) for the restriction
of \(u_n\) to \({}_{\ell}T'_n\). For fixed \(\ell\) and varying \(n\), the morphisms
\({}_{\ell}u_n\) form a projective system and therefore arise from a unique
\(S\)-group morphism
\[
{}_{\ell}u:{}_{\ell}T'\longrightarrow H
\]
(\SGARefLink{sga3:xv:proposition:1-6}{1.6} a)). Since \(H\) is separated (Exposé VI\textsubscript{B}, \SGARefLink{prop:sga3-VIB-5.2}{5.2}) and \(u_0\) is
a monomorphism, \({}_{\ell}u\) is a monomorphism (\SGARefLink{sga3:ix:theorem:6:8}{Exposé IX, 6.8}), and even a closed
immersion because it is finite, \({}_{\ell}T'\) being finite over \(S\). Denote its
image by \(M(\ell)\). The family of subgroups of multiplicative type \(M(\ell)\) is
clearly coherent in the sense of \SGARefLink{sga3:xv:theorem:4-1}{4.1}. Put
\[
C_\ell=\operatorname{Centr}_H(M(\ell)).
\]
This is represented by a closed subgroup prescheme (\SGARefLink{sga3:xv:lemma:2:5}{2.5}), since \(H\) is separated,
and is smooth over \(S\) (\SGARefLink{sga3:xi:corollary:2:4}{Exposé XI, 2.4}). The \(C_\ell\) form a filtered decreasing
family of closed subschemes and are therefore stationary, since \(H\) is noetherian.
Their stationary value is a smooth closed subgroup \(C\) such that
\(C_s=\operatorname{Centr}_{H_s}(T_s)\), by the density theorem (\SGARefLink{sga3:ix:theorem:4:7}{Exposé IX, 4.7}).
It remains to show that for every point \(t\) of \(S\), \(C_t\) is the centralizer
in \(H_t\) of a subtorus \(T_t\); this will also imply that \(C_t\) is connected.
This follows from the more precise lemma below, applied to the family
\(M(\ell)_t\) of subgroups of \(H_t\).

\SGARefTarget{sga3:xv:lemma:6:8}
\textbf{Lemma 6.8.} \emph{Let \(G\) be a connected algebraic group over a field
\(k\), let \(r>0\) be an integer, and let \(q\) be an integer prime to the
characteristic\footnote{Editor's note: the word ``residual'' has been removed.}
of \(k\). Let \(M(\ell)\), where \(\ell\) ranges over the powers of \(q\), be a
coherent family of subgroups of \(G\), of multiplicative type and of type
\((\mathbb Z/\ell\mathbb Z)^r\) (cf.~\SGARefLink{sga3:xv:lemma:4:6}{4.6}). Let \(M\) be the algebraic subgroup of
\(G\) generated by the \(M(\ell)\) (\emph{loc.~cit.}), and let \(T\) be the unique
maximal torus of \(M\) (cf.~\SGARefLink{sga3:xv:lemma:3:4}{3.4}). Then, for all sufficiently large \(\ell\),
\[
\operatorname{Centr}_G(T)=\operatorname{Centr}_G(M)
  =\operatorname{Centr}_G(M(\ell)).
\]}


The last equality is clear. To prove the first, introduce the centre \(Z\) of
\(G\), put \(G'=G/Z\), let \(M'\) (resp. \(T'\)) be the image of \(M\) (resp.
\(T\)) in \(G'\), and let \(K\) be the inverse image of \(T'\) in \(G\), that is,
the algebraic subgroup of \(G\) generated by \(T\) and \(Z\). It plainly suffices
to prove \(K\supset M\), or equivalently \(T'=M'\). Now \(M\) is smooth and
connected (\SGARefLink{sga3:xv:lemma:4:6}{4.6}), while \(G'\) is affine (\SGARefLink{sga3:xii:6.1}{Exposé XII, 6.1}); hence \(M'\) is the
direct product of its maximal torus \(T'\) (\SGARefLink{sga3:xii:6.6}{Exposé XII, 6.6} d)) and a unipotent
group (\emph{Bible} 4, Th.~4), after assuming \(k\) algebraically closed. The image
of \(M(\ell)\) in \(M'\) is therefore contained in \(T'\). Thus the inverse image
of \(T'\) in \(M\) contains \(M(\ell)\) for every \(\ell\), and consequently is
equal to \(M\). Hence \(M'=T'\). This proves \SGARefLink{sga3:xv:lemma:6:8}{6.8}, and therefore \SGARefLink{sga3:xv:lemma:6:7}{6.7}.

We now prove \SGARefLink{sga3:xv:lemma:6:6}{6.6}. We have reduced to the case where \(S\) is the spectrum of a
discrete valuation ring \(A\), which we may moreover assume complete with
algebraically closed residue field. After replacing \(A\), if necessary, by its
normalization in a finite extension of its fraction field, we may assume that
\((G_t)_{\mathrm{red}}\) is smooth.\footnote{Editor's note: details or references
should be supplied here.} To prove \SGARefLink{sga3:xv:lemma:6:6}{6.6} we may clearly replace \(G\) by the identity
component of the schematic closure in \(G\) of \((G_t)^0_{\mathrm{red}}\), and
therefore assume that \(G\) is flat over \(S\), has connected fibres, and has smooth
generic fibre \(G_t\).

\begin{enumerate}
\renewcommand{\labelenumi}{\alph{enumi})}
\item Suppose that \(H_s\) is the centralizer in \((G_s)_{\mathrm{red}}\) of a torus
\(T_s\). We show that \(H_t\) is then the centralizer in \(G_t\) of a subtorus of
\(G_t\). By \SGARefLink{sga3:xv:lemma:6:7}{Lemma 6.7} there is a subgroup scheme \(C\) of \(H\), smooth over \(S\),
whose closed fibre is \(H_s\), such that
\(C_t=\operatorname{Centr}_{H_t}(T_t)\) for a subtorus \(T_t\) of \(H_t\). Since
\(H\) is smooth over \(S\) with connected fibres, dimension considerations give
\(C=H\). With the notation of \SGARefLink{sga3:xv:lemma:6:7}{6.7}, \(H=C=C_\ell\) for all sufficiently large
\(\ell\). Similarly, put
\(C'_\ell=\operatorname{Centr}_G(M(\ell))\) (\SGARefLink{sga3:xv:lemma:2:5}{2.5}), and let \(C'\) be the stationary
value of the \(C'_\ell\) for large \(\ell\) (\SGARefLink{sga3:xv:lemma:2-5-bis}{2.5 bis}). The group scheme \(C'\)
contains \(H\), and
\[
C'_t=\operatorname{Centr}_{G_t}(T_t)\quad\text{(6.8)},\qquad
C'_s=\operatorname{Centr}_{G_s}(T_s).
\]
The hypothesis on \(H_s\) gives \(\dim H_s=\dim C'_s\). Moreover,
\(\dim H_s=\dim H_t\), since \(H\) is smooth over \(S\), and
\(\dim C'_t\leq \dim C'_s\) (Exposé VI\textsubscript{B}, \SGARefLink{prop:sga3-VIB-4.1}{4.1}). Hence
\(\dim H_t=\dim C'_t\). Since \(G_t\) is smooth and connected,
\(C'_t=\operatorname{Centr}_{G_t}(T_t)\) is smooth and connected; finally,
\[
H_t=C'_t=\operatorname{Centr}_{G_t}(T_t).
\]

\item Suppose that \(H_s\) contains a Cartan subgroup \(C_s\) of \(G_s\), that is,
the centralizer in \((G_s)_{\mathrm{red}}\) of a maximal torus \(T_s\) of \(G_s\).
By \SGARefLink{sga3:xv:lemma:6:7}{6.7} there is a subgroup scheme \(C\) of \(H\), smooth over \(S\) with connected
fibres, lifting \(C_s\). It suffices to prove that \(C_t\) contains a Cartan subgroup
of \(G_t\). By part a), applied with \(H=C\), \(C_t\) is the centralizer in \(G_t\)
of a subtorus of \(G_t\), and hence contains a Cartan subgroup of \(G_t\).

\item Suppose that \(H_s\) is a parabolic subgroup of \(G_s\). Let
\(N=\operatorname{Norm}_G(H)\). By hypothesis \(H\) is equal to its connected
normalizer, so \(N_s\) is smooth and consequently
\[
N_s=\operatorname{Norm}_{(G_s)_{\mathrm{red}}}(H_s)=H_s
\qquad\text{(\SGARefLink{sga3:xii:7.10}{Exposé XII, 8 bis}).}
\]
Thus \(N\) is flat over \(S\). We shall see in Exposé XVI that under these
conditions \(G/N\) is representable. Since \(H_s=N_s\) is a parabolic subgroup,
\((G/N)_s\) is proper. Since \(G/N\) is flat over \(S\) with connected fibres,
EGA III, 5.5.1 implies that \(G/N\) is proper over \(S\). Thus
\((G/N)_t=G_t/N_t\) is proper over \(t\), and so is \(G_t/H_t\), because
\(N_t/H_t\) is finite. It remains only to apply the following lemma.
\end{enumerate}

\SGARefTarget{sga3:xv:lemma:6:9}
\textbf{Lemma 6.9.} \emph{Let \(k\) be a field, \(G\) a \(k\)-algebraic group,
\(H\) a smooth connected algebraic subgroup of \(G\), and
\(N=\operatorname{Norm}_G(H)\). If \(\dim H=\dim N\) and \(G/H\) is proper, then
\(H\) is a parabolic subgroup of \(G\).}

Indeed, we may assume that \(k\) is algebraically closed and that \(G\) is smooth
and connected. The centre \(Z\) of \(G\) is contained in \(N\), and
\(\dim H=\dim N\) implies that \(Z'=(Z)^0_{\mathrm{red}}\) is contained in \(H\).
Hence \(G'=G/Z'\) is affine (\SGARefLink{sga3:xii:6.1}{Exposé XII, 6.1}). Replacing \(G\) by \(G'\) and \(H\)
by its image in \(G'\), we reduce to the case in which \(G\) is affine
(\SGARefLink{sga3:xiv:4.9}{Exposé XIV, 4.9}), and \SGARefLink{sga3:xv:lemma:6:9}{Lemma 6.9} then follows from \emph{Bible} 6, Th.~4.
This proves \SGARefLink{sga3:xv:lemma:6:6}{Lemma 6.6}.


To complete the proof of \SGARefLink{sga3:xv:theorem:6:4}{6.4}, we must show that the \(S\)-preschemes representing
\(\mathcal L_C\) (respectively \(\mathcal{CT}\) and \(\mathcal P\)) are of finite
presentation over \(S\). This assertion is local on \(S\), so we may assume first
that \(S\) is affine and then, since \(G\) is of finite presentation, that \(S\) is
noetherian (EGA IV, 8.9). We have just seen that the natural inclusions
\(\mathcal{CT}\to\mathcal L\) and \(\mathcal P\to\mathcal L\) are immersions.
The same is therefore true of
\(\mathcal{CT}\to\mathcal L_C\) and \(\mathcal P\to\mathcal L_C\).
Consequently, it is enough to prove that \(\mathcal L_C\) is represented by an
\(S\)-prescheme of finite presentation.

Resume the notation introduced in \SGARefLink{sga3:xv:theorem:5-2}{5.2}. For every integer \(n\geq 0\), let
\(\mathcal L^n\) be the subfunctor of \(\mathcal L\) defined by
\[
\mathcal L^n(S')
 =\left\{H\in\mathcal L(S')\ \middle|\ 
   \operatorname{Norm}_{G_{S'}}(H)
   =\operatorname{Norm}_{G_{S'}}\bigl(H^{(n)}\bigr)\right\}.
\]
The \(S\)-functor \(\mathcal L^n\) is thus represented by an open subprescheme
of \(\mathcal L\), the disjoint union of the \(\mathcal L_r^n\).
Each \(\mathcal L_r^n\) is of finite presentation over \(S\) (\SGARefLink{sga3:xv:lemma:5:3}{5.3}), and is empty
for
\[
r>\sup_{s\in S}\dim G_s,
\]
which is a finite number because \(S\) is quasi-compact. Hence \(\mathcal L^n\)
is of finite presentation over \(S\). It remains to prove that
\(\mathcal L_C\) is contained in \(\mathcal L^n\) for all sufficiently large
\(n\).

For each point \(s\) of \(S\), let \(d(s)\) be the least integer \(n\), finite or
infinite, such that
\[
\mathcal L_{C,G_s}\subset\mathcal L^n_{G_s}.
\]
It is enough to show that \(d\) is bounded on \(S\). Indeed, if \(M\) is an
upper bound, then \(\mathcal L_C\) is set-theoretically contained in
\(\mathcal L^M\), and hence is contained in \(\mathcal L^M\), since the latter
is open in \(\mathcal L\). A direct constructibility argument reduces us to
showing that, when \(S\) is noetherian and integral with generic point \(t\),
the function \(d\) is bounded on a neighbourhood of \(t\).

\begin{enumerate}
\renewcommand{\labelenumi}{\alph{enumi})}
\item \textbf{Reduction to the case in which \(G\) is smooth over \(S\).}
Proceeding as in \SGARefLink{sga3:xv:proposition:6-2}{6.2}, we see that, after possibly changing \(S\), we may assume
that \(G_{\mathrm{red}}\) is a group prescheme smooth over \(S\); denote it by
\(G'\). Put
\[
X=\mathcal L_{C,G},\qquad X'=\mathcal L_{C,G'},
\]
and let \(H\) (respectively \(H'\)) be the subgroup prescheme of \(G_X\)
(respectively \(G'_{X'}\)) universal for the functor
\(\mathcal L_{C,G}\) (respectively \(\mathcal L_{C,G'}\)).
Since \(H\) is smooth over \(X\), the prescheme
\(H\times_X X_{\mathrm{red}}\) is reduced, and is therefore contained in
\(G'_{X_{\mathrm{red}}}\). It is an element of
\(\mathcal L_{C,G'}(X_{\mathrm{red}})\), and hence defines a canonical
morphism
\[
p:X_{\mathrm{red}}\longrightarrow X'.
\]
It is clear that \(p\) is a monomorphism; we show that it is in fact an
immersion. Let \(N'\) be the normalizer of \(H'\) in \(G_{X'}\).
The set of points \(s\in X'\) for which
\(H'_s\to N'_s\) is an open immersion is an open subset \(U\), and
\[
H'|_U\longrightarrow N'|_U
\]
is an open immersion (Exposé VI\textsubscript{B}, 2.5, and EGA IV, 17.9.5).
It follows from 5.1(iii) that \(U\) is the largest open subset of \(X'\)
over which \(H'\) equals its connected normalizer in \(G_{X'}\).
Thus \(H'|_U\in\mathcal L_{C,G}(U)\), and it follows at once that \(p\)
is an isomorphism of \(X_{\mathrm{red}}\) onto \(U_{\mathrm{red}}\).
If \(\mathcal L_{C,G}\) is known to be of finite type when \(G\) is smooth,
then \(X'\) is of finite type over \(S\); consequently
\(X_{\mathrm{red}}\) is of finite type (since \(S\) is noetherian), and is
therefore contained in \(\mathcal L_G^M\) for all sufficiently large \(M\).
The same then holds for \(X\).

\item \textbf{The case in which \(S\) is the spectrum of an algebraically
closed field \(k\) of characteristic \(p>0\), and \(G\) is a smooth
algebraic group over \(k\).}
Instead of the infinitesimal neighbourhoods \(H^{(n)}\) of a subgroup
prescheme \(H\) of \(G\), we use the radicial subgroups \(F_n(H)\), the
kernels of the iterates of the Frobenius morphism of \(H\)
(Exposé VII\textsubscript{A}, 4). This is legitimate here because \(F_n(H)\)
is contained in the infinitesimal neighbourhood of order \(p^n\) of the
unit section of \(H\). If \(T\) is a subtorus of \(G\), then
\[
F_n(T)={}_{p^n}T,
\]
and duality immediately shows that
\(\operatorname{Centr}_G({}_{p^n}T)=\operatorname{Centr}_G(T)\) for all
sufficiently large \(n\). We have the following more precise proposition.
\end{enumerate}

\SGARefTarget{sga3:xv:proposition:6:10}
\textbf{Proposition 6.10.} \emph{Let \(k\) be a field of characteristic
\(p>0\), let \(G\) be a smooth \(k\)-algebraic group, let \(T\) be a maximal
torus of \(G_{\bar k}\), and let \(m\) be the least integer such that
\[
\operatorname{Centr}_G({}_{p^m}T)^0
 =\operatorname{Centr}_G(T)^0.
\]
Then, for every \(k\)-prescheme \(S\) and every
\(H\in\mathcal L_C(S)\), one has
\[
\operatorname{Norm}_{G_S}(H)
 =\operatorname{Norm}_{G_S}\bigl(F_m(H)\bigr).
\]
In particular, \(\mathcal L_C\) is contained in \(\mathcal L^{p^m}\).}

Since \(H\) is smooth and of finite presentation over \(S\), and has constant
nilpotent rank, namely that of \(G\), we shall see in the next section (\SGARefLink{sga3:xv:theorem:7-3}{7.3})
that the functor \(\mathcal C_H\) of Cartan subgroups of \(H\) is represented
by an \(S\)-prescheme smooth over \(S\). The reader will verify that the proof
given there does not use the finite-type property of \(\mathcal L_{C,H}\).
It follows from \SGARefLink{sga3:xiii:theorem:3:1}{Exposé XIII, 3.1}, that we may consider the open subset \(U\)
of regular points of \(H\).

Let \(S'\) be an \(S\)-prescheme and let \(g\in G(S')\) normalize
\(F_m(H)_{S'}\). To prove that \(g\) normalizes \(H_{S'}\), it is enough to
prove that
\[
\operatorname{int}(g)U_{S'}\subset H_{S'}.
\]
Indeed, \(H_{S'}\cap\operatorname{int}(g)H_{S'}\) would then contain an open
subgroup of \(\operatorname{int}(g)H_{S'}\), and hence would equal it, since
the latter has connected fibres. After replacing \(S'\) by a suitable
covering \(S''\), and then replacing \(S''\) by \(S\), we are reduced to
proving the following: if \(g\in G(S)\) normalizes \(F_m(H)\) and
\(u\in U(S)\), then \(\operatorname{int}(g)u\in H(S)\).

Let \(C\) be the unique Cartan subgroup of \(H\) ``containing'' \(u\)
(\SGARefLink{sga3:xiii:corollary:3:2}{Exposé XIII, 3.2}). It is enough to show that
\[
C'=\operatorname{int}(g)C\subset H.
\]
Since \(H\in\mathcal L_{C,G}(S)\), the group \(C\) is also a Cartan subgroup
of \(G\). The latter admits maximal tori, so \(C\) is the centralizer in
\(G_S^0\) of its unique maximal torus \(T\) (\SGARefLink{sga3:xii:7.1:item:a}{Exposé XII, 7.1(a)}).
By the definition of \(m\), and because any two maximal tori of \(G\) are
locally conjugate for the fpqc topology (\SGARefLink{sga3:xii:7.1}{Exposé XII, 7.1}),
\[
C=\operatorname{Centr}_G(T)^0
 =\operatorname{Centr}_G({}_{p^m}T)^0.
\]
Conjugating by \(g\), we obtain
\[
C'=\operatorname{Centr}_G\bigl(\operatorname{int}(g)({}_{p^m}T)\bigr)^0.
\]
Now \(\operatorname{int}(g)({}_{p^m}T)\) is a subgroup of multiplicative
type of
\[
\operatorname{int}(g)\bigl(F_m(H)\bigr)=F_m(H),
\]
because \(g\) normalizes \(F_m(H)\); it is therefore contained in \(H\).
\SGARefLink{sga3:xiii:theorem:2:1}{Exposé XIII, 2.1} (proved there when the base is a field, but extending at
once to an arbitrary base) then shows that, in order to prove
\(C'\subset H\), it is enough to prove
\[
\operatorname{Lie}C'\subset\operatorname{Lie}H.
\]
But
\[
\operatorname{Lie}C'
 =\operatorname{int}(g)(\operatorname{Lie}C)
 \subset\operatorname{int}(g)(\operatorname{Lie}H).
\]
On the other hand, we may assume \(m\geq1\), and Exposé
VII\textsubscript{A}, \SGARefLink{sga3:VIIA:4.1.2}{4.1.2}, gives
\[
\operatorname{Lie}H
 =\operatorname{Lie}\bigl(F_m(H)\bigr)
 =\operatorname{Lie}\bigl(\operatorname{int}(g)(F_m(H))\bigr)
 =\operatorname{int}(g)(\operatorname{Lie}H).
\]
Thus \(\operatorname{Lie}C'\subset\operatorname{Lie}H\), which completes
the proof of \SGARefLink{sga3:xv:proposition:6:10}{6.10}.


We shall need another description of the integer \(m\) introduced in \SGARefLink{sga3:xv:proposition:6:10}{6.10}.

\SGARefTarget{sga3:xv:lemma:6:11}
\textbf{Lemma 6.11.} \emph{Let \(G\) be a smooth algebraic group over an
algebraically closed field \(k\) of characteristic \(p>0\), let \(T\) be a
maximal torus of \(G\), put \(\mathfrak g=\operatorname{Lie}G\), and let \(R\)
be the set of nonzero characters of \(T\) occurring in the representation of
\(T\) on \(\mathfrak g\) induced by the adjoint representation of \(G\).
For \(r\in R\), let \(e_r\) be the greatest integer \(n\) such that \(p^n\)
divides \(r\) in the character group of \(T\). Then
\[
m=\sup_{r\in R}(e_r+1)
\]
if \(R\neq\varnothing\), and \(m=0\) otherwise.}

Indeed, \(\operatorname{Centr}_G(T)\) is smooth and is contained in
\(\operatorname{Centr}_G({}_{p^m}T)\). Hence
\[
\begin{aligned}
\operatorname{Centr}_G(T)^0
 &=\operatorname{Centr}_G({}_{p^m}T)^0\\
&\Longleftrightarrow\
\operatorname{Lie}\bigl(\operatorname{Centr}_G(T)\bigr)
 =\operatorname{Lie}\bigl(\operatorname{Centr}_G(F_m(T))\bigr)\\
&\Longleftrightarrow\
\mathfrak g^T=\mathfrak g^{{}_{p^m}T}
\qquad\text{(Exposé II, 5.2.3).}
\end{aligned}
\]
With the usual notation,
\[
\mathfrak g=\mathfrak g_0\oplus\bigoplus_{r\in R}\mathfrak g_r.
\]
Thus \(\mathfrak g^T=\mathfrak g_0\), whereas
\[
\mathfrak g^{{}_{p^m}T}
 =\mathfrak g_0\oplus\bigoplus_{r\in R'}\mathfrak g_r,
\]
where \(R'\) is the subset of \(R\) formed by the characters whose restriction
to \({}_{p^m}T\) is zero. A nonzero character \(r\) of \(T\) has trivial
restriction to \({}_{p^m}T\) if and only if \(m\leq e_r\), which proves the
lemma.

\begin{enumerate}
\renewcommand{\labelenumi}{\alph{enumi})}
\setcounter{enumi}{2}
\item \textbf{Return to the proof of 6.4.}
By part a) and the paragraph preceding it, we have reduced to the case in
which \(S\) is a noetherian integral scheme and \(G\) is smooth over \(S\).
We must show that \(d\) is bounded on a neighbourhood of the generic point
\(t\) of \(S\). After possibly changing \(S\), we may assume that \(G\) has
a split maximal torus (\SGARefLink{sga3:xv:proposition:6-2}{6.2})
\[
T=D_S(M).
\]
Let
\[
\mathfrak g=\bigoplus_{\lambda\in M}\mathfrak g^\lambda
\]
be the decomposition of the Lie algebra of \(G\) according to the characters
of \(T\), and let \(R\) be the finite set of nonzero \(\lambda\in M\) for
which \(\mathfrak g^\lambda\neq0\). We distinguish two cases.
\end{enumerate}

\textbf{First case: \(t\), and hence every point of \(S\), has positive
residue characteristic \(p\).}
It is clear from the foregoing that \(d\) is then bounded by \(p^m\), where
\(m\) is defined as in \SGARefLink{sga3:xv:lemma:6:11}{6.11}.

\textbf{Second case: \(t\) has residue characteristic zero.}
For each \(\lambda\in R\), let \(n_\lambda\) be the greatest integer dividing
\(\lambda\) in the group \(M\), and put
\[
n=\prod_{\lambda\in R}n_\lambda.
\]
For every prime number \(q\) dividing \(n\), let \(S_q\) be the closed subset
of \(S\) formed by the points of residue characteristic \(q\), and let \(U\)
be the nonempty open subset, containing \(t\), complementary to the union of
the \(S_q\). If \(s\in U\), then either \(s\) has residue characteristic zero,
in which case \(d(s)=1\) by \SGARefLink{sga3:xv:remark:5:6}{5.6}, or \(s\) has residue characteristic
\(p>0\) not dividing \(n\). In the latter case the integer \(m\) attached to
\(G_s\) in \SGARefLink{sga3:xv:lemma:6:11}{6.11} is at most one. Moreover, Exposé
VII\footnote{Editor's note: the argument should be made explicit.} gives
\[
\operatorname{Norm}_{G_{S'}}\bigl(F_1(H)\bigr)
 =\operatorname{Norm}_{G_{S'}}(\operatorname{Lie}H)
 =\operatorname{Norm}_{G_{S'}}\bigl(H^{(1)}\bigr).
\]
Finally, \SGARefLink{sga3:xv:proposition:6:10}{6.10} shows that, for
\(H\in\mathcal L_{C,G_s}(S')\),
\[
\operatorname{Norm}_{G_{S'}}(H)
 =\operatorname{Norm}_{G_{S'}}\bigl(H^{(1)}\bigr).
\]
Consequently \(d(s)\leq1\), so \(d\) is bounded by \(1\) on \(U\).
This completes the proof of \SGARefLink{sga3:xv:theorem:6:4}{6.4}.

\SGARefTarget{sga3:xv:corollary:6:12}
\textbf{Corollary 6.12.} \emph{Let \(S\) be a prescheme and \(G\) an
\(S\)-prescheme in groups of finite presentation. Suppose that the nilpotent
rank of the fibres of \(G\) (respectively the dimension of their Borel
subgroups) is locally constant on \(S\). Then the functor \(\mathcal C\) of
Cartan subgroups of \(G\) (respectively the functor \(\mathcal B\) of Borel
subgroups of \(G\)) is represented by an \(S\)-prescheme of finite
presentation over \(S\).}

Indeed, after restricting \(S\), we may assume that the nilpotent rank
\(\nu\) of the fibres is constant. Then \(\mathcal C\) is represented by the
open-and-closed subprescheme of the prescheme \(X\) representing
\(\mathcal L_C\) (\SGARefLink{sga3:xv:theorem:6:4}{6.4}) over which the universal subgroup of \(G_X\) for the
functor \(\mathcal L_C\) has relative dimension \(\nu\). The proof is
analogous for \(\mathcal B\), using the representability of \(\mathcal P\).
